% Bagian: computability
% Bab: recursive-functions
% Subbagian: halting-problem

\documentclass[../../../include/open-logic-section]{subfiles}

\begin{document}

\olfileid[id]{cmp}{rec}{hlt}
\olsection{Masalah Penghentian}

Secara umum, \emph{masalah penghentian} ialah masalah memutuskan, apabila
diberikan spesifikasi~$e$ (misalnya, program) dari suatu prosedur komputasi
yang mungkin hanya mendefinisikan fungsi parsial dan suatu bilangan~$n$, apakah
komputasi itu pada masukan~$n$ berhenti, yakni menghasilkan suatu keluaran.
Alan Turing membuktikan hasil terkenal bahwa masalah ini sendiri tidak dapat
diselesaikan oleh suatu prosedur yang dapat dihitung;
dengan kata lain, fungsi
\[
h(e, n) =
\begin{cases}
1 & \text{jika komputasi $e$ berhenti pada masukan $n$}\\
0 & \text{selain itu}
\end{cases}
\]
tidak dapat dihitung.

Dalam konteks fungsi rekursif parsial, peran spesifikasi program dapat
dimainkan oleh indeks~$e$ yang diberikan dalam teorema bentuk normal Kleene.
Jika $f$ merupakan fungsi rekursif parsial, setiap $e$ yang membuat persamaan
dalam teorema bentuk normal berlaku merupakan suatu indeks dari~$f$. Jika
diberikan suatu bilangan~$e$, teorema bentuk normal menyatakan bahwa
\[
\cfind{e}(x) \simeq U(\mu s \; T(e, x, s))
\]
bersifat rekursif parsial, dan bagi setiap fungsi rekursif parsial
$f\colon\Nat\to\Nat$, terdapat suatu $e\in\Nat$ sedemikian sehingga
$\cfind{e}(x)\simeq f(x)$ bagi semua~$x\in\Nat$. Bahkan, bagi setiap~$f$
semacam itu, bukan hanya ada satu, melainkan tak hingga banyak~$e$. Fungsi
\emph{penghentian}~$h$ didefinisikan oleh
\[
h(e, x) =
\begin{cases}
1 & \text{jika $\cfind{e}(x) \fdefined$}\\
0 & \text{selain itu.}
\end{cases}
\]
Perhatikan bahwa $h(e,x)=0$ jika dan hanya jika
$\cfind{e}(x)\fundefined$.

\begin{thm}
\ollabel{thm:halting-problem}
Fungsi penghentian~$h$ tidak bersifat rekursif parsial.
\end{thm}

\begin{proof}
Andaikan $h$ bersifat rekursif parsial. Kita dapat mendefinisikan
\[
d(y) =
\begin{cases}
1 & \text{jika $h(y, y) = 0$}\\
\umin{x}{x \neq x} & \text{selain itu.}
\end{cases}
\]
Karena tidak ada bilangan $x$ yang memenuhi $x\neq x$, tidak ada
$\umin{x}{x\neq x}$, sehingga $d(y)\fundefined$ jika dan hanya jika
$h(y,y)\neq0$. Dari definisi ini diperoleh
\begin{enumerate}
\item $d(y)\fdefined$ jika dan hanya jika
  $\cfind{y}(y)\fundefined$.
\item $d(y)\fundefined$ jika dan hanya jika
  $\cfind{y}(y)\fdefined$.
\end{enumerate}
Jika $h$ bersifat rekursif parsial, maka $d$ juga bersifat rekursif parsial.
Jadi, menurut teorema bentuk normal Kleene, $d$ mempunyai suatu indeks~$e_d$.
Tinjau nilai $h(e_d,e_d)$. Ada dua kemungkinan, yaitu $0$ dan~$1$.
\begin{enumerate}
\item Jika $h(e_d,e_d)=1$, maka $\cfind{e_d}(e_d)\fdefined$. Namun,
  $\cfind{e_d}\simeq d$, sedangkan $d(e_d)$ terdefinisi jika dan hanya jika
  $h(e_d,e_d)=0$. Ini suatu kontradiksi.
\item Jika $h(e_d,e_d)=0$, maka $\cfind{e_d}(e_d)\fundefined$. Namun,
  $\cfind{e_d}\simeq d$, sedangkan $d(e_d)$ terdefinisi jika dan hanya jika
  $h(e_d,e_d)=0$. Ini juga suatu kontradiksi.
\end{enumerate}
Kedua kemungkinan menghasilkan kontradiksi. Maka, $h$ tidak mungkin bersifat
rekursif parsial.
\end{proof}

\end{document}
